% part: intuitionistic-logic
% chapter: introduction
% section: natural-deduction

\documentclass[../../../include/open-logic-section]{subfiles}

\begin{document}

\olfileid{int}{int}{ntd}

\olsection{नैसर्गिक निगमन}

$\FalseCl$ चे नियम वगळलेले नैसर्गिक निगमन ही
अंतःप्रज्ञावादी तर्कशास्त्रासाठीची प्रमाणित
!!{derivation} व्यवस्था आहे. येथे ते नियम पुन्हा
देऊन BHK अर्थनिर्धारणाच्या साहाय्याने त्यामागची
प्रेरणा स्पष्ट करू. प्रत्येक नियम असा समजू शकतो
की आधारविधानांच्या रचना असतील, तर निष्कर्षाचीही
रचना आहे असे त्यावरून ठरवता येते.

नैसर्गिक निगमनातील !!{derivation}s मध्ये मुक्त न
केलेली गृहीतके असतात. म्हणून~$\Gamma$ या
!!{undischarged} गृहीतकांपासून~$!A$ ची
!!a{derivation} असेल, तर तिला प्रत्येक $!B \in
\Gamma$ च्या रचना घेऊन~$!A$ ची रचना देणारे
फलन मानावे. कोणत्याही !!{undischarged} गृहीतकांशिवाय
$!A$ ची !!a{derivation} असेल, तर BHK अर्थनिर्धारणाच्या
अर्थाने~$!A$ ची रचना असते. मात्र पुढील चर्चेत गरज
नसेल तेव्हा~$\Gamma$ चा उल्लेख वगळू.

गृहीतक~$!A$ स्वतःच त्या !!{undischarged}
गृहीतक~$!A$ पासून~$!A$ ची !!a{derivation} आहे.
हे BHK अर्थनिर्धारणाशी जुळते: रचनांवरील ओळख फलन
$!A$ ची कोणतीही रचना घेऊन~$!A$ चीच रचना देते.

\subsection{संयोग}

\begin{defish}
\AxiomC{$!A$}
\AxiomC{$!B$}
\RightLabel{\Intro{\land}}
\BinaryInfC{$!A \land !B$}
\DisplayProof
\hfill
\begin{tabular}{l}
\AxiomC{$!A \land !B$}
\RightLabel{\Elim{\land}}
\UnaryInfC{$!A$}
\DisplayProof
\\[3ex]
\AxiomC{$!A \land !B$}
\RightLabel{\Elim{\land}}
\UnaryInfC{$!B$}
\DisplayProof
\end{tabular}
\end{defish}

आपल्याकडे अनुक्रमे $!A_1$ आणि $!A_2$ यांच्या
रचना $N_1$, $N_2$ आहेत असे समजा. मग
$!A_1 \land !A_2$ चीही रचना आहे: जोडी
$\tuple{N_1, N_2}$.

BHK अर्थनिर्धारणानुसार $!A_1 \land !A_2$ ची
रचना म्हणजे $\tuple{N_1, N_2}$ ही जोडी. अशी जोडी
आपल्याकडे आहे असे समजा. मग प्रत्येक संयुक्त घटकाची
रचना आहे: $N_1$ ही~$!A_1$ ची आणि $N_2$ ही
$!A_2$ ची रचना आहे.

\subsection{सशर्त}

\begin{defish}
  \AxiomC{$\Discharge{!A}{u}$}
  \DeduceC{$!B$}
  \DischargeRule{$\Intro{\lif}$}{u}
  \UnaryInfC{$!A \lif !B$}
  \DisplayProof
\hfill
  \AxiomC{$!A \lif !B$}
  \AxiomC{$!A$}
  \RightLabel{$\Elim{\lif}$}
  \BinaryInfC{$!B$}
  \DisplayProof
\end{defish}

!!{undischarged} गृहीतक~$!A$ पासून~$!B$ ची
!!a{derivation} असेल, तर~$!A$ च्या रचनांना~$!B$
च्या रचनांत नेणारे फलन~$f$ असते. तेच फलन
$!A \lif !B$ ची रचना आहे. म्हणून $\Intro\lif$
च्या आधारविधानाची रचना~$!A$ च्या रचनेवर अवलंबून
असेल, तर निष्कर्ष $!A
\lif !B$ ची रचना असते.

दुसरीकडे, $!A$ ची रचना~$N$ आणि $!A \lif !B$
ची रचना~$f$ आहेत असे समजा. $!A \lif !B$ ची
रचना हे~$!A$ च्या रचनांना~$!B$ च्या रचनांत नेणारे
फलन असते. म्हणून $f(N)$ ही~$!B$ ची रचना आहे;
म्हणजे $\Elim\lif$ च्या निष्कर्षाची रचना आहे.

\subsection{विकल्पयोग}

\begin{defish}
\begin{tabular}{l}
\AxiomC{$!A$}
\RightLabel{\Intro{\lor}}
\UnaryInfC{$!A \lor !B$}
\DisplayProof
\\[3ex]
\AxiomC{$!B$}
\RightLabel{\Intro{\lor}}
\UnaryInfC{$!A \lor !B$}
\DisplayProof
\end{tabular}
\hfill
\AxiomC{$!A \lor !B$}
\AxiomC{$\Discharge{!A}{n}$}
\DeduceC{$!C$}
\AxiomC{$\Discharge{!B}{n}$}
\DeduceC{$!C$}
\DischargeRule{\Elim{\lor}}{n}
\TrinaryInfC{$!C$}
\DisplayProof
\end{defish}

$!A_i$ ची रचना~$N_i$ असेल, तर तिच्यापासून
$!A_1 \lor !A_2$ ची रचना~$\tuple{i, N_i}$ करता
येते. उलट, $!A_1 \lor !A_2$ ची रचना म्हणजे
$\tuple{i, N_i}$ ही जोडी आपल्याकडे असेल, जिथे
$N_i$ ही~$!A_i$ ची रचना आहे; तसेच अनुक्रमे
$!A_1$, $!A_2$ यांच्या रचनांना~$!C$ च्या रचनांत
नेणारी फलने $f_1$, $f_2$ असतील, तर $f_i(N_i)$
ही~$!C$ ची रचना मिळते. हाच~$\Elim\lor$ चा निष्कर्ष.

\subsection{असंगती}

\begin{defish}
  \AxiomC{$\lfalse$}
  \RightLabel{\FalseInt}
  \UnaryInfC{$!A$}
  \DisplayProof
\end{defish}

!!{undischarged} गृहीतके~$!B_1$, \dots, $!B_n$
यांपासून~$\lfalse$ ची !!a{derivation} असेल, तर
$!B_1$, \dots, $!B_n$ यांच्या रचनांना~$\lfalse$
च्या रचनेत नेणारे फलन~$f(M_1,
\dots, M_n)$ असते. $\lfalse$ ची रचना नसल्याने
$!B_1$, \dots, $!B_n$ या सर्वांच्या रचनाही एकत्र
असू शकत नाहीत. म्हणून $f$ बाबत पुढीलही खरे असते:
$M_1$, \dots, $M_n$ या अनुक्रमे $!B_1$, \dots,
$!B_n$ यांच्या रचना \emph{असतील तर}, $f(M_1, \dots,
M_n)$ ही~$!A$ ची रचना \emph{असते}.

\subsection{$\lnot$ साठीचे नियम}

$\lnot !A$ ची व्याख्या $!A \lif \lfalse$ अशी
असल्याने काटेकोरपणे पाहता~$\lnot$ साठी वेगळे
नियम लागत नाहीत. तरी ते द्यायचे झाले, तर असे दिसतील:

\begin{defish}
\AxiomC{$\Discharge{!A}{n}$}
\noLine
\DeduceC{$\lfalse$}
\DischargeRule{\Intro{\lnot}}{n}
\UnaryInfC{$\lnot !A$}
\DisplayProof
\hfill
\AxiomC{$\lnot !A$}
\AxiomC{$!A$}
\RightLabel{\Elim{\lnot}}
\BinaryInfC{$\lfalse$}
\DisplayProof
\end{defish}


\subsection{\usetoken{P}{derivation} ची उदाहरणे}

\begin{enumerate}
\item $\Proves !A \lif (\lnot !A \lif \lfalse)$, i.e., $\Proves !A
  \lif ((!A \lif \lfalse) \lif \lfalse)$
  \begin{prooftree}
    \AxiomC{$\Discharge{!A}{2}$}
    \AxiomC{$\Discharge{!A \lif \lfalse}{1}$}
    \RightLabel{\Elim\lif}
    \BinaryInfC{$\lfalse$}
    \DischargeRule{\Intro\lif}{1}
    \UnaryInfC{$(!A \lif \lfalse)\lif \lfalse$}
    \DischargeRule{\Intro\lif}{2}
    \UnaryInfC{$!A \lif (!A \lif \lfalse) \lif \lfalse$}
  \end{prooftree}

\item $\Proves ((!A \land !B) \lif !C) \lif (!A \lif (!B \lif !C))$
  \begin{prooftree}
    \AxiomC{$\Discharge{(!A \land !B) \lif !C}{3}$}
    \AxiomC{$\Discharge{!A}{2}$}
    \AxiomC{$\Discharge{!B}{1}$}
    \RightLabel{\Intro\land}
    \BinaryInfC{$!A \land !B$}
    \RightLabel{\Elim\lif}
    \BinaryInfC{$!C$}
    \DischargeRule{\Intro\lif}{1}
    \UnaryInfC{$!B \lif !C$}
    \DischargeRule{\Intro\lif}{2}
    \UnaryInfC{$!A \lif (!B \lif !C)$}
    \DischargeRule{\Intro\lif}{3}
    \UnaryInfC{$((!A \land !B) \lif !C) \lif (!A \lif (!B \lif !C))$}
  \end{prooftree}

\item $\Proves \lnot(!A \land \lnot !A)$, i.e., $\Proves (!A \land (!A
  \lif \lfalse))\lif \lfalse$
  \begin{prooftree}
    \AxiomC{$\Discharge{!A \land (!A \lif \lfalse)}{1}$}
    \RightLabel{\Elim\land}
    \UnaryInfC{$!A \lif \lfalse$}
    \AxiomC{$\Discharge{!A \land (!A \lif \lfalse)}{1}$}
    \RightLabel{\Elim\land}
    \UnaryInfC{$!A$}
    \RightLabel{\Elim\lif}
    \BinaryInfC{$\lfalse$}
    \DischargeRule{\Intro\lif}{1}
    \UnaryInfC{$(!A \land (!A \lif \lfalse))\lif \lfalse$}
  \end{prooftree}

\item $\Proves \lnot\lnot(!A \lor \lnot !A)$, i.e., $\Proves ((!A \lor
  (!A \lif \lfalse))\lif \lfalse)\lif \lfalse$
  \begin{prooftree}\hskip-3em
    \AxiomC{$\Discharge{(!A \lor (!A \lif \lfalse))\lif \lfalse}{2}$}
    \AxiomC{$\Discharge{(!A \lor (!A \lif \lfalse))\lif \lfalse}{2}$}
    \AxiomC{$\Discharge{!A}{1}$}
    \RightLabel{\Intro\lor}
    \UnaryInfC{$!A \lor (!A \lif \lfalse)$}
    \RightLabel{\Elim\lif}
    \BinaryInfC{$\lfalse$}
    \DischargeRule{\Intro\lif}{1}
    \UnaryInfC{$!A \lif \lfalse$}
    \RightLabel{\Intro\lor}
    \UnaryInfC{$!A \lor (!A \lif \lfalse)$}
    \RightLabel{\Elim\lif}
    \insertBetweenHyps{\hskip -4em}
    \BinaryInfC{$\lfalse$}
    \DischargeRule{\Intro\lif}{2}
    \UnaryInfC{$((!A \lor (!A \lif \lfalse))\lif \lfalse)\lif \lfalse$}
  \end{prooftree}
\end{enumerate}

\begin{prop}
  अंतःप्रज्ञावादी तर्कशास्त्रातील नैसर्गिक निगमनात
  $\Gamma \Proves !A$ असेल, तर अभिजात तर्कशास्त्रातही
  $\Gamma \Proves !A$ असते. विशेषतः, $!A$ हे
  अंतःप्रज्ञावादी प्रमेय असेल, तर ते अभिजात प्रमेयही असते.
\end{prop}

\begin{proof}
  नैसर्गिक निगमनाचा प्रत्येक नियम अभिजात नैसर्गिक
  निगमनाचाही नियम आहे. म्हणून अंतःप्रज्ञावादी
  तर्कशास्त्रातील प्रत्येक !!{derivation} ही अभिजात
  तर्कशास्त्रातीलही !!a{derivation} आहे.
\end{proof}

\begin{prob}
  पुढील !!{formula}s ची अंतःप्रज्ञावादी तर्कशास्त्रात
  !!{derivation}s द्या:
  \begin{enumerate}
    \item $(\lnot !A \lor !B) \lif (!A \lif !B)$
    \item $\lnot\lnot\lnot !A \lif \lnot !A$
    \item $\lnot\lnot (!A \land !B) \liff (\lnot\lnot !A\land \lnot \lnot !B)$
    \item $\lnot (!A \lor !B) \liff (\lnot !A \land \lnot !B)$
    \item $(\lnot !A \lor \lnot !B) \lif \lnot (!A \land !B)$
    \item $\lnot\lnot ( !A \land !B) \lif  (\lnot\lnot !A \lor
    \lnot\lnot !B)$
  \end{enumerate}
\end{prob}

\end{document}
